
%%%%%%%%%%%%%%%%%%%%%%%%%%%%%%%%%%%%%%%%%%%%%%%%%%%%%%%%%%%%%
%% Begin exercise %%
%%%%%%%%%%%%%%%%%%%%%%%%%%%%%%%%%%%%%%%%%%%%%%%%%%%%%%%%%%%%%
\ex{Power IGBTs}

%%%%%%%%%%%%%%%%%%%%%%%%%%%%%%%%%%%%%%%%%%%%%%%%%%%%%%%%%%%%%
%% Task 1 %%
%%%%%%%%%%%%%%%%%%%%%%%%%%%%%%%%%%%%%%%%%%%%%%%%%%%%%%%%%%%%%

\task{When is the IGBT better than the MOSFET?}

Two devices are available for a converter leg.

\begin{table}[h]
  \centering
  % Erste Tabelle
  \begin{minipage}{0.45\textwidth}
    \centering
    \begin{tabular}{ll}
        \toprule
        \multicolumn{2}{l}{\textbf{MOSFET}} \\
        \midrule
        Drain-source blocking voltage $V_\mathrm{DSS}$ & $\SI{600}{\volt}$\\
        On-state resistance $R_\mathrm{DS(on)}$: & $\SI{80}{\milli\ohm}$ \\
        Switching time $(t_\mathrm{r} + t_\mathrm{f})$ & $\SI{80}{\nano\second}$\\
        \multicolumn{2}{l}{}  \\        
        \bottomrule
    \end{tabular}
  \end{minipage}
  \hfill % Abstand zwischen den Tabellen
  % Zweite Tabelle
  \begin{minipage}{0.45\textwidth}
    \centering
    \begin{tabular}{ll}
        \toprule
        \multicolumn{2}{l}{\textbf{IGBT}}  \\
        \midrule
        Maximal CE voltage $V_\mathrm{CE,max}$: & $\SI{1200}{\volt}$\\
        On-state resistance $V_\mathrm{CE(sat)}$: $\SI{2.1}{\volt}$ \\
        Switching energy $E_\mathrm{on} + E_\mathrm{off}$: & $\SI{1.8}{\milli\joule}$ \\
        \multicolumn{2}{l}{(Switching energy at $\SI{400}{\volt},\SI{20}{\ampere}$)}  \\        
        \bottomrule
    \end{tabular}
  \end{minipage}
\caption{MOSFET and IGBT data.}
\label{table:mosfet_igbt_data}        
\end{table}

\textbf{Application I}: $V_\mathrm{DC} = \SI{400}{\volt}$, $I = \SI{20}{\ampere}$, $f_\mathrm{s} = \SI{80}{\kilo\hertz}$, duty ratio $D = 0.5$ \\
\textbf{Application II}: $V_\mathrm{DC} = \SI{800}{\volt}$, $I = \SI{50}{\ampere}$, $f_\mathrm{s} = \SI{5}{\kilo\hertz}$ \\

For Application II, compare the IGBT with a hypothetical high-voltage MOSFET having
\begin{equation*}
    R_\mathrm{DS(on)} = \SI{180}{\milli\ohm}
\end{equation*}
and assume the IGBT switching energy is $\SI{6}{\milli\joule}$ per cycle.

\subtask{For Application I, compare approximate losses and choose the preferred device.}

\begin{solutionblock}
    The MOSFET conduction loss is
    \begin{align}
        I_\mathrm{rms} &= \SI{20}{\ampere}\sqrt{0.5} = \SI{14.14}{\ampere} \\
        P_\mathrm{ON,M} &= \left( \SI{14.14}{\ampere} \right)^2 \cdot \SI{80}{\milli\ohm} \approx \SI{16}{\watt} \\
        P_\mathrm{sw,M} &= \frac{1}{6} \cdot \SI{400}{\volt} \cdot \SI{20}{\ampere} \cdot \SI{80}{\nano\second} \cdot \SI{80}{\kilo\hertz}
        \approx \SI{8.53}{\watt} \\
        P_\mathrm{tot,M} &\approx \SI{24.5}{\watt}.
    \end{align}

    For the IGBT:
    \begin{align}
        P_\mathrm{ON,I} &= \SI{2.1}{\volt} \cdot \SI{20}{\ampere} \cdot 0.5 = \SI{21}{\watt} \\
        P_\mathrm{sw,I} &= \SI{1.8}{\milli\joule} \cdot \SI{80}{\kilo\hertz} = \SI{144}{\watt} \\
        P_\mathrm{tot,I} &\approx \SI{165}{\watt}.
    \end{align}

    Therefore, at \SI{400}{\volt} and \SI{80}{\kilo\hertz}, the MOSFET wins strongly.
\end{solutionblock}

\subtask{For Application II, compare approximate losses and choose the preferred device.}

\begin{solutionblock}
    For the hypothetical high-voltage MOSFET:
    \begin{align}
        I_\mathrm{rms} &= \SI{50}{\ampere}\sqrt{0.5} = \SI{35.36}{\ampere} \\
        P_\mathrm{ON,M} &= \left( \SI{35.36}{\ampere} \right)^2 \cdot \SI{180}{\milli\ohm} \approx \SI{225}{\watt}.
    \end{align}

    For the IGBT, taking a realistic high-current value of $V_\mathrm{CE(sat)} \approx \SI{2.4}{\volt}$:
    \begin{align}
        P_\mathrm{ON,I} &= \SI{2.4}{\volt} \cdot \SI{50}{\ampere} \cdot 0.5 = \SI{60}{\watt} \\
        P_\mathrm{sw,I} &= \SI{6}{\milli\joule} \cdot \SI{5}{\kilo\hertz} = \SI{30}{\watt} \\
        P_\mathrm{tot,I} &\approx \SI{90}{\watt}.
    \end{align}

    Hence, at high voltage and low switching frequency, the IGBT is clearly preferred.
\end{solutionblock}

\subtask{Explain physically why the preferred device changes with voltage and frequency.}

\begin{solutionblock}
    The IGBT gains lower conduction loss at high voltage because conductivity modulation reduces the effective resistance of the drift region.
    However, the stored minority charge produces larger switching losses, especially at high frequency.
    Therefore, MOSFETs dominate at lower voltage and high switching frequency, while IGBTs become attractive at high voltage and lower frequency.
\end{solutionblock}

\subtask{State the main design insight students should retain.}

\begin{solutionblock}
    The preferred device cannot be decided from one parameter alone. 
    Voltage class, frequency, and the balance between conduction and switching loss must all be considered together.
\end{solutionblock}


%%%%%%%%%%%%%%%%%%%%%%%%%%%%%%%%%%%%%%%%%%%%%%%%%%%%%%%%%%%%%
%% Task 2 %%
%%%%%%%%%%%%%%%%%%%%%%%%%%%%%%%%%%%%%%%%%%%%%%%%%%%%%%%%%%%%%

\task{Conductivity modulation made quantitative}

A \SI{1200}{\volt} IGBT carries $I_\mathrm{C} = \SI{40}{\ampere}$. The corresponding parameters are listed in \autoref{table:igbt_task2_device}.


\begin{table}[ht]
    \centering
    \begin{tabular}{ll}
        \toprule
        \multicolumn{2}{l}{\textbf{IGBT}}  \\
        \midrule
        Junction voltage $V_\mathrm{J1}$: & $\SI{0.95}{\volt}$ \\
        Drift region voltage $V_\mathrm{drift}$: & $\SI{0.55}{\volt}$ \\
        Channel resistance $R_\mathrm{channel}$: & $\SI{12}{\milli\ohm}$ \\
        \bottomrule
    \end{tabular}
    \caption{IGBT semiconductor characteristic.}
    \label{table:igbt_task2_device}
\end{table}

Use the simplified on-state model $V_\mathrm{CE(on)} \approx V_\mathrm{J1} + V_\mathrm{drift} + I_\mathrm{C} R_\mathrm{channel}$. \\
A comparable high-voltage MOSFET would have $R_\mathrm{DS(on)} = \SI{90}{\milli\ohm}$. \\

\subtask{Estimate $V_\mathrm{CE(on)}$ for the IGBT.}

\begin{solutionblock}
    \begin{equation}
        V_\mathrm{CE(on)} = \SI{0.95}{\volt} + \SI{0.55}{\volt} + \SI{40}{\ampere} \cdot \SI{12}{\milli\ohm}
        = \SI{1.98}{\volt}.
    \end{equation}
\end{solutionblock}

\subtask{Estimate the on-state drop of the MOSFET at $\SI{40}{\ampere}$.}

\begin{solutionblock}
    \begin{equation}
        V_\mathrm{DS(on)} = I_\mathrm{D} R_\mathrm{DS(on)}
        = \SI{40}{\ampere} \cdot \SI{90}{\milli\ohm}
        = \SI{3.6}{\volt}.
    \end{equation}
\end{solutionblock}

\subtask{For duty ratio $D = 0.5$, compare the conduction losses.}

\begin{solutionblock}
    \begin{align}
        P_\mathrm{cond,IGBT} &= \SI{1.98}{\volt} \cdot \SI{40}{\ampere} \cdot 0.5 = \SI{39.6}{\watt} \\
        P_\mathrm{cond,MOSFET} &= \SI{3.6}{\volt} \cdot \SI{40}{\ampere} \cdot 0.5 = \SI{72}{\watt}.
    \end{align}

    Thus, the IGBT roughly halves the conduction loss in this case.
\end{solutionblock}

\subtask{Explain in words what conductivity modulation is buying the designer.}

\begin{solutionblock}
    Injected holes and electrons raise the carrier concentration in the drift region, thereby reducing its effective resistance.
    This is why the IGBT becomes attractive when the drift region would otherwise dominate the MOSFET on-state loss.
\end{solutionblock}


%%%%%%%%%%%%%%%%%%%%%%%%%%%%%%%%%%%%%%%%%%%%%%%%%%%%%%%%%%%%%
%% Task 3 %%
%%%%%%%%%%%%%%%%%%%%%%%%%%%%%%%%%%%%%%%%%%%%%%%%%%%%%%%%%%%%%

\task{Latch-up is not a mystery; it is a voltage-drop problem}

A student says: "Latch-up sounds like an abstract failure mechanism."
Show that it can be framed as a simple lateral voltage-drop problem.\\

In an IGBT cell, assume the hole current in the p-body is $I_\mathrm{h}$, and the effective body-spreading resistance is $R_\mathrm{b}$.
Latch-up risk begins when the lateral voltage drop is large enough to forward-bias the $n^+$-source / p-body junction, approximately $\SI{0.7}{\volt}$.

\begin{table}[ht]
    \centering
    \begin{tabular}{lll}
        \toprule
        Parameter  & \textbf{Case A} & \textbf{Case B} \\
        \midrule
        Hole current $I_\mathrm{h}$: & $\SI{12}{\ampere}$ & $\SI{12}{\ampere}$  \\
        Effective body-spreading resistance $R_\mathrm{b}$: & $\SI{60}{\milli\ohm}$ & $\SI{35}{\milli\ohm}$ \\
        \bottomrule
    \end{tabular}
    \caption{IGBT semiconductor characteristic.}
    \label{table:igbt_task2_device}
\end{table}

\subtask{Compute the lateral p-body voltage drop in each case.}

\begin{solutionblock}
    For Case A:
    \begin{equation}
        V_\mathrm{b,A} = I_\mathrm{h} R_\mathrm{b}
        = \SI{12}{\ampere} \cdot \SI{60}{\milli\ohm}
        = \SI{0.72}{\volt}.
    \end{equation}

    For Case B:
    \begin{equation}
        V_\mathrm{b,B} = I_\mathrm{h} R_\mathrm{b}
        = \SI{12}{\ampere} \cdot \SI{35}{\milli\ohm}
        = \SI{0.42}{\volt}.
    \end{equation}
\end{solutionblock}

\subtask{Which case is vulnerable to latch-up, and why?}

\begin{solutionblock}
    Case A is vulnerable because \SI{0.72}{\volt} is enough to forward-bias the source-body junction.
    Case B is much safer because the lateral voltage drop stays below the approximate forward-bias threshold.
\end{solutionblock}

\subtask{Suggest two structural ways to increase latch-up immunity.}

\begin{solutionblock}
    Latch-up immunity can be improved by reducing the body-spreading resistance, keeping the source width small, and providing an easier hole bypass path.
    Any such measure keeps the lateral voltage drop below the forward-bias threshold of the parasitic npn path.
\end{solutionblock}

\subtask{Explain why gate control is lost once the parasitic thyristor turns on.}

\begin{solutionblock}
    Once the coupled pnp and npn sections both conduct, the parasitic thyristor structure is triggered.
    In that condition, the current no longer depends on the MOS gate channel, so the gate loses effective control over the device current.
\end{solutionblock}


%%%%%%%%%%%%%%%%%%%%%%%%%%%%%%%%%%%%%%%%%%%%%%%%%%%%%%%%%%%%%
%% Task 4 %%
%%%%%%%%%%%%%%%%%%%%%%%%%%%%%%%%%%%%%%%%%%%%%%%%%%%%%%%%%%%%%

\task{PT or NPT? Selecting a device structure}

Two IGBTs are available.

\begin{table}[ht]
    \centering
    \begin{tabular}{lll}
        \toprule
        Parameter  & \textbf{Device P (PT-IGBT)} & \textbf{Device N (NPT-IGBT)} \\
        \midrule
        Switch off energy loss $E_\mathrm{off}$: & $\SI{2.5}{\milli\joule}$  & $\SI{4}{\milli\joule}$  \\
        Tail-current fraction: & $\SI{15}{\percent}$ & $\SI{25}{\percent}$ \\
        Short-circuit withstand time: & $\SI{6}{\micro\second}$ & $\SI{10}{\micro\second}$ \\
        Avalanche ruggedness : & moderate & high \\
        \bottomrule
    \end{tabular}
    \caption{IGBT semiconductor characteristic.}
    \label{table:igbt_task4_device}
\end{table}

Assume both have similar $E_\mathrm{on} = \SI{1.5}{\milli\joule}$.\\

\textbf{Application I}: Servo inverter with $\SI{600}{\volt}$ , $\SI{25}{\ampere}$, $f_\mathrm{s} = \SI{20}{\kilo\hertz}$ \\
\textbf{Application II}: Industrial motor drive  with $\SI{1200}{\volt}$ , $\SI{40}{\ampere}$, $f_\mathrm{s} = \SI{4}{\kilo\hertz}$ \\
Remark: For Application II overloads and harsh fault conditions are expected.

\subtask{Compute the switching-loss power in each application for each device.}

\begin{solutionblock}
    Switching power is
    \begin{equation}
        P_\mathrm{sw} = f_\mathrm{s} (E_\mathrm{on} + E_\mathrm{off}).
    \end{equation}

    For Device P:
    \begin{equation}
        E_\mathrm{sw,P} = \SI{1.5}{\milli\joule} + \SI{2.5}{\milli\joule} = \SI{4.0}{\milli\joule}.
    \end{equation}

    For Device N:
    \begin{equation}
        E_\mathrm{sw,N} = \SI{1.5}{\milli\joule} + \SI{4.0}{\milli\joule} = \SI{5.5}{\milli\joule}.
    \end{equation}

    Application I:
    \begin{align}
        P_\mathrm{sw,P} &= \SI{4.0}{\milli\joule} \cdot \SI{20}{\kilo\hertz} = \SI{80}{\watt} \\
        P_\mathrm{sw,N} &= \SI{5.5}{\milli\joule} \cdot \SI{20}{\kilo\hertz} = \SI{110}{\watt}.
    \end{align}

    Application II:
    \begin{align}
        P_\mathrm{sw,P} &= \SI{4.0}{\milli\joule} \cdot \SI{4}{\kilo\hertz} = \SI{16}{\watt} \\
        P_\mathrm{sw,N} &= \SI{5.5}{\milli\joule} \cdot \SI{4}{\kilo\hertz} = \SI{22}{\watt}.
    \end{align}
\end{solutionblock}

\subtask{Choose the preferred device for each application.}

\begin{solutionblock}
    At \SI{20}{\kilo\hertz}, Device P (PT-IGBT) is preferred because its switching-loss advantage is substantial.

    At \SI{4}{\kilo\hertz}, the switching-loss penalty of Device N (NPT-IGBT) is small, while its ruggedness advantage matters much more.
    Hence, Application II should use the NPT device.
\end{solutionblock}

\subtask{Explain the structural reason for the choice and why the PT/NPT choice is really a ``speed versus ruggedness'' trade-off.}

\begin{solutionblock}
    PT structures use a buffer layer and generally tend toward faster switching and shorter tail current.
    NPT structures tend toward better ruggedness, avalanche capability, and thermal stability.

    Therefore, PT versus NPT is fundamentally a trade-off between switching speed and fault robustness.
\end{solutionblock}


%%%%%%%%%%%%%%%%%%%%%%%%%%%%%%%%%%%%%%%%%%%%%%%%%%%%%%%%%%%%%
%% Task 5 %%
%%%%%%%%%%%%%%%%%%%%%%%%%%%%%%%%%%%%%%%%%%%%%%%%%%%%%%%%%%%%%

\task{Device-limit and SOA screening}

For an IGBT the simplified limits are listed in \autoref{table:igbt_task5_limits}.
Use the following simplified IGBT limits:


\begin{table}[ht]
    \centering
    \begin{tabular}{ll}
        \toprule
        Parameter  & Value \\
        \midrule
        collector-emitter blocking voltage $V_\mathrm{CES}$: & $\SI{1200}{\volt}$ \\
        maximal gate-emitter voltage $\lvert V_\mathrm{GE(max)} \rvert$: & $\SI{20}{\volt}$ \\
        continuous collector current limit $I_\mathrm{C}$: & $\SI{75}{\ampere}$\\
        pulsed current limit $I_\mathrm{CM}$: & $\SI{150}{\ampere}$ \\
        short-circuit withstand time $t_\mathrm{SC}$: & $\SI{8}{\micro\second}$ \\
        dc conduction thermal proxy: & $V_\mathrm{CE} I_\mathrm{C} \leq \SI{250}{\watt}$ \\
        $\SI{100}{\micro\second}$ pulse limit: & $V_\mathrm{CE} I_\mathrm{C} \leq \SI{4000}{\watt}$ \\
        turn-off RBSOA proxy for $V_\mathrm{CE} > \SI{400}{\volt}$: & $I_\mathrm{C} \leq \frac{1200}{V_\mathrm{CE}}$ \\
        \bottomrule
    \end{tabular}
    \caption{IGBT semiconductor limits.}
    \label{table:igbt_task5_limits}
\end{table}

Consider the following cases:
\begin{itemize}
    \setlength{\itemsep}{-2pt} % Distance between the points (Part1)
    \setlength{\parskip}{-2pt} % Distance between the points (Part2)    
    \item \textbf{A.} $V_\mathrm{CE} = \SI{2.2}{\volt}$, $I_\mathrm{C} = \SI{60}{\ampere}$, dc
    \item \textbf{B.} $V_\mathrm{CE} = \SI{600}{\volt}$, $I_\mathrm{C} = \SI{3}{\ampere}$, turn-off event, \SI{100}{\micro\second}
    \item \textbf{C.} $V_\mathrm{CE} = \SI{600}{\volt}$, $I_\mathrm{C} = \SI{1.5}{\ampere}$, turn-off event, \SI{100}{\micro\second}
    \item \textbf{D.} short-circuit event, detection trips only after \SI{12}{\micro\second}
\end{itemize}

\subtask{Decide which cases are safe.}

\begin{solutionblock}
    Case A:
    \begin{equation}
        P = \SI{2.2}{\volt} \cdot \SI{60}{\ampere} = \SI{132}{\watt} < \SI{250}{\watt}.
    \end{equation}
    Therefore, Case A is safe.

    Case B:
    \begin{equation}
        P = \SI{600}{\volt} \cdot \SI{3}{\ampere} = \SI{1800}{\watt} < \SI{4000}{\watt}.
    \end{equation}
    But for RBSOA:
    \begin{equation}
        I_\mathrm{max} = \frac{1200}{600} = \SI{2}{\ampere}.
    \end{equation}
    Since the actual current is \SI{3}{\ampere}, Case B is unsafe.

    Case C:
    \begin{equation}
        P = \SI{600}{\volt} \cdot \SI{1.5}{\ampere} = \SI{900}{\watt} < \SI{4000}{\watt}.
    \end{equation}
    The RBSOA current limit remains \SI{2}{\ampere}, and the actual current is only \SI{1.5}{\ampere}.
    Therefore, Case C is safe.

    Case D is unsafe because the short circuit is not cleared until \SI{12}{\micro\second}, while the device is rated only for \SI{8}{\micro\second}.
\end{solutionblock}

\subtask{Identify which violated mechanism dominates in each unsafe case, and explain why RBSOA is more restrictive than simple thermal power during turn-off.}

\begin{solutionblock}
    In Case B, the dominant violated mechanism is the reverse-bias safe operating area. 
    In Case D, the dominant violated mechanism is the short-circuit withstand limit.

    RBSOA is more restrictive than simple thermal power during turn-off because large $\mathrm{d}v_\mathrm{CE}/\mathrm{d}t$, stored carriers, and latch-up-related effects appear during the switching transient. 
    Thus, a point may look thermally acceptable and still be dynamically unsafe.
\end{solutionblock}

\subtask{Explain why short-circuit withstand time is a real design parameter, not just a datasheet decoration.}

\begin{solutionblock}
    The short-circuit withstand time tells the designer how long the protection system may need to detect and clear a fault without destroying the device.
    Therefore, it is a real coordination parameter between the device, the gate driver, and the fault-protection strategy.
\end{solutionblock}
